\documentclass[conference]{IEEEtran}
\usepackage[utf8]{inputenc}
\usepackage[T1]{fontenc}
\usepackage{cite}
\usepackage{amsmath,amssymb}
\usepackage{graphicx}
\usepackage{booktabs}
\usepackage{url}

\title{Measuring Validator and Sandboxed-Execution Coverage for LLM‑Generated NetOps Artifacts: A Paired Confirmatory Study}

\author{
\IEEEauthorblockN{Autonomous AI Researcher}
\IEEEauthorblockA{CNSM 2026 Agentic AI Researcher Track}
}

\begin{document}

\maketitle

\begin{abstract}
We preregistered and executed a paired confirmatory experiment (study\_id f2c3d8b1-7a6e-4a9a-b2c3-5f8e1d2b6c7e) that measured the marginal detection benefit of adding sandboxed emulation to a deterministic validator for LLM-generated intent\ensuremath{\rightarrow}configuration artifacts. Under shared\_initial\_candidate semantics using the declared model "gpt-5-mini" (provider: openai\_responses) we evaluated Npairs = 720 confirmatory tasks. The deterministic validator (deterministic\_netops\_validator\_v5) flagged 719/720 = 0.9986111111111111; the guarded pipeline (validator + sandboxed emulation) flagged 720/720 = 1.0. Paired difference (guarded - baseline) = 0.001388888888888884; exact McNemar two-sided p = 1.0; 95\% bootstrap percentile CI (10000 resamples, seed = 1729) = [0.0, 0.004166666666666667]. Run accounting: model\_calls\_used = 721 (720 shared initial + 1 repair-stage call). These artifact-grounded results lie far below our preregistered minimum effect-of-interest (0.20); we report the preregistered design, exact quantitative outcomes, run-accounting diagnostics, artifact-grounded interpretation for the negligible guarded benefit in this regime, and reproducibility pointers into the archived execution bundle.
\end{abstract}

\section{Introduction}
Generative large language models (LLMs) are increasingly proposed to translate operator intent into device configuration and NetOps workflows. Operational deployment requires generated artifacts to satisfy syntactic correctness, workflow ordering, and higher-order safety/policy constraints. A common mitigation architecture is generate\ensuremath{\rightarrow}validate\ensuremath{\rightarrow}(repair)\ensuremath{\rightarrow}execute, sometimes augmented by sandboxed emulation to reveal runtime-only failures that static validators miss. Prior surveys and system reports advocate such workflow-centred mitigations [1]--[3], but provide limited large-scale paired measures of the marginal value added by sandboxed execution. Motivated by this gap, we preregistered and executed a controlled paired study to measure the aggregate detection-rate difference contributed uniquely by an automated sandboxed-emulation stage when the same initial generated candidate is analyzed by both conditions.

Research question and contributions. Our preregistered research question was: for synthetic intent\ensuremath{\rightarrow}configuration tasks under shared\_initial\_candidate semantics, what is the paired detection-rate difference (guarded - baseline) between (A) a deterministic validator-only pipeline and (B) the same deterministic validator followed by automated sandboxed emulation? Contributions: (1) a machine-readable preregistration and faithful autonomous execution; (2) an artifact-grounded paired analysis on Npairs = 720 with complete accounting; (3) an artifact-grounded interpretation explaining a near-zero marginal effect; (4) concise reproducibility pointers into the archived bundle and prescriptive boundary conditions where guarded stages are likely to matter.

\section{Related Work}
Workflow-centred mitigations and verification tools. Recent survey and architectural reports document the practical trend of combining LLM generation with deterministic validators, emulators, and repair loops to raise operational assurance for NetOps artifacts [1]. These works emphasize that correctness in intent\ensuremath{\rightarrow}configuration pipelines is multidimensional: syntactic parsing, command-ordering/workflow constraints, invariants (e.g., keep‑one‑uplink‑up), and dynamic runtime properties may each require different V\&V primitives. Tool-focused case studies highlight the evolution and practical limits of static verifiers: mature configuration-analysis systems have progressively expanded rule-sets and topology-aware checks but still exhibit gaps for runtime-dependent behaviors and stateful sequences [2]. That evolution shows why researchers and operators often layer an emulator or sandbox after static validation: emulation can exercise temporal behaviors, device state transitions, and interactions that are difficult to capture exhaustively as static rules.

Empirical and prototype evidence. Prototype systems that integrate generation with verification and execution-harnesses report heterogeneous gains: in some task regimes a verifier alone catches most dangerous errors, while in others emulation uncovers subtle runtime failures or environment-specific side effects that static rules miss [1], [3]. The practical implication drawn across this literature is that the marginal utility of a guarded (validator+emulation) pipeline depends on (a) validator coverage relative to the task distribution, (b) the fidelity of the emulator, and (c) the distributional properties of the input intents (e.g., constrained, parser-validated intents versus noisy natural-language requests). These prior empirical and system-level observations motivated a paired experimental design: by reusing the identical initial generated candidate across conditions (shared\_initial\_candidate semantics) we isolate only downstream guarded-stage effects (emulation observations or permitted repair invocations) rather than conflating them with generation variability.

Gaps and evaluation needs. While surveys and system reports argue for generate\ensuremath{\rightarrow}validate\ensuremath{\rightarrow}repair and for emulator-assisted validation as best practice for high-assurance NetOps, they do not provide large-scale, preregistered paired estimates of the marginal detection benefit under controlled generation semantics. The literature also documents important open questions that our study was designed to probe: how much additional coverage does sandboxed emulation provide when a strong deterministic validator already targets the same workflow patterns, and which task regimes (adversarial inputs, cross‑vendor CLI heterogeneity, or stateful multi-step procedures) are most likely to reveal nontrivial guarded-stage benefits? By situating our paired confirmatory design amid these prior findings, we provide an artifact-grounded quantitative datapoint addressing one narrow slice of the broader evaluation agenda advocated in [1]--[3].

\section{Methodology}
Preregistration and analytic intent. We preregistered a confirmatory paired-binary design (study\_id f2c3d8b1-7a6e-4a9a-b2c3-5f8e1d2b6c7e). The primary estimand was the paired\_success\_rate\_difference\_guarded\_minus\_baseline aggregated across confirmatory samples; the preregistered minimum effect-of-interest for power planning was 0.20 absolute. The preregistration (frozen prior to execution) specified: inclusion/exclusion rules, deterministic missingness and reseed policies (up to two automatic reseeds for parser failures), repair-stage limits (<=1 repair call per task), contamination detection and deterministic replacement rules, and the analysis executor (paired\_binary\_analysis\_v1) with paired inference procedures and pre-registered sensitivity analyses.

Task bank, topology templates, and vendor abstraction. Confirmatory items were synthetic intent\ensuremath{\rightarrow}configuration tasks generated from the intent\_configuration\_repair\_v1 family. The bank covered three abstracted vendor‑CLI families and three topology templates (leaf--spine, 3-node service chain, simple triangle). Each task manifest encoded: the initial device state, an explicit intent, a required command sequence, allowed-touched fields, and parser-acceptance constraints. Inclusion required vendor-specific parser acceptance after up to two deterministic reseeds; irrecoverable parser failures were to be deterministically excluded and replaced per the preregistered missingness plan. In the executed run there were zero irrecoverable parser exclusions and complete\_pair\_count = 720 (analysis/missingness\_summary.csv).

Model interface, sampling semantics, and shared\_initial\_candidate enforcement. The declared/requested model was recorded as "gpt-5-mini" and the execution model\_configuration.json contains declared\_model\_version = "gpt-5-mini", max\_output\_tokens = 2000, and temperature = null (sha256 a40e96e212ab87da...). Important operational clarification (preregistered and enforced): temperature = null in the archived configuration indicates the client did not set an explicit temperature parameter; null does not imply temperature=0 or provider-side determinism. To isolate guarded-stage effects we enforced shared\_initial\_candidate semantics: for each pair the pipeline performed a single sampled initial generation (shared) which both the baseline validator and the guarded pipeline consumed. Concretely, provider\_call\_audit and deterministic\_reconciliation artifacts record stage\_counts: shared\_initial\_generation = 720 and repair = 1, and model\_calls\_used = 721 (720 shared + one repair-stage call). Because the same sampled candidate was analyzed by both conditions for each pair, any between-condition discordance must stem from (i) guarded-stage sandboxed emulation observations not encoded as validator violations, and/or (ii) the permitted repair-stage intervention recorded in run accounting.

Validator implementation and scoring rule. Baseline validation used deterministic\_netops\_validator\_v5 (scorer\_version = 5.0) executing the preregistered scoring rule full\_intent\_constraint\_validation. The validator performs multi-level checks recorded in per-episode validator\_trace objects: vendor-specific parsing, syntactic command parsing, command ordering/workflow checks against required\_sequence, required-field presence, dependency constraints (e.g., must-shutdown-before-mtu-change patterns), and policy constraints (protected interfaces, group-wise minimum-active constraints). The validator emits structured outputs used in scoring: normalized\_configuration, parsed\_command\_count, required\_command\_count, observed\_touched\_field\_count, violation\_count, and violation\_codes. Per-episode artifacts (execution/scoring/task-XXXXX-*.json) show these fields and validator\_version = 5.0 for all confirmatory episodes.

Sandboxed-emulation harness and guarded pipeline semantics. The guarded pipeline applied the baseline deterministic validator and then, for candidates not blocked by parser failure, executed an automated sandboxed-emulation harness that instantiated an emulated representation of the declared topology and applied the candidate commands. The emulation harness produced runtime observations that were normalized and summarized into the same validator-trace structure (validation\_before/validation\_after, final\_state, parsed\_command\_count, observed\_touched\_field\_count). The emulation stage was designed to detect runtime-only symptoms such as transient ordering-dependent state changes, reachability/regression signals, or state transitions that static rules do not explicitly capture. The preregistered repair policy permitted up to one repair-stage model call per task; repair calls (if used) and their provider traces are recorded in the execution manifest and deterministic\_reconciliation artifacts.

Inclusion/exclusion, missingness, and contamination controls. The preregistered missingness plan allowed two automatic deterministic reseeds for parser failures; irrecoverable parser failures would be deterministically replaced. Contamination controls included deterministic SHA-256 artifact hashing, automated prompt-template divergence checks, and validator/emulator binary/config checksum verification before confirmatory execution. The execution manifest and analysis/contamination\_summary.csv report no contamination flags and no irrecoverable parser exclusions for the confirmatory set.

Statistical analysis procedures. The primary confirmatory analysis used the paired\_binary\_analysis\_v1 executor specified in the preregistration. The primary test computed the aggregate paired detection-rate difference (guarded - baseline) across complete pairs and reported an exact McNemar two-sided p-value. Pre-registered uncertainty quantification used paired-bootstrap percentile confidence intervals with 10,000 resamples and a fixed seed (bootstrap\_seed = 1729) for reproducibility. Secondary exploratory analyses (per-class and per-vendor paired differences) were pre-registered with Bonferroni adjustments; these are archived but not the focus of the confirmatory report. The preregistered default treatment of irrecoverable technical failures for the primary analysis was to exclude complete-pair technical losses (complete\_pair\_primary); a sensitivity analysis treats irrecoverable failures as nondetections (documented in the preregistration). The executed analysis followed the preregistered complete\_pair\_primary plan; analysis artifacts and logs are archived (analysis/analysis\_log.jsonl and analysis/paired\_contingency\_table.csv).

Deterministic reconciliation and run accounting. To enable artifact-grounded claims and independent checking we recorded provider\_call\_audit and deterministic\_reconciliation artifacts. These reconcile: planned\_episode\_count = 1440, completed\_episode\_count = 1440, complete\_pair\_count = 720, failed\_episode\_count = 0, model\_calls\_used = 721, and stage\_counts (shared\_initial\_generation = 720; repair = 1). These machine-verifiable run-accounting artifacts are referenced throughout the manuscript when we state numerical outcomes and when tracing the single discordant guarded-only detection back to a documented repair-stage invocation.

\section{Results}
Primary confirmatory counts and test statistics (artifact-grounded). The preregistered paired analysis produced the following exact contingency counts (analysis/paired\_contingency\_table.csv, sha256 efe6e8e7\allowbreak{}016fa843\allowbreak{}a8226e91\allowbreak{}1dbde829\allowbreak{}e943903d\allowbreak{}719101da\allowbreak{}8d956b3f\allowbreak{}f899133c): n11 = 719 (detected by both), n10 = 1 (guarded-only), n01 = 0 (baseline-only), n00 = 0 (neither). Marginal detection rates are: baseline = 719/720 = 0.9986111111111111 and guarded = 720/720 = 1.0. The paired estimate (guarded - baseline) = 0.001388888888888884. The exact McNemar two-sided test reported p = 1.0. Paired-bootstrap percentile 95\% CI (10000 resamples, bootstrap\_seed = 1729) = [0.0, 0.004166666666666667]. All numeric values above are reproduced verbatim from analysis/results.json and analysis artifacts.

Bootstrap and test details. The bootstrap CI was computed with 10000 resamples (seed = 1729) using paired resampling of episode-pairs; the analysis executor used a paired-bootstrap-percentile approach (artifact: analysis/analysis\_log.jsonl). The McNemar exact two-sided p-value complements the bootstrap CI by testing symmetry in the discordant-pair counts; here the observed discordant counts (n10 = 1, n01 = 0) produce no evidence of a systematic guarded advantage at conventional significance thresholds (p = 1.0 under the exact test reported by the executor).

Run accounting and discordant-pair provenance. Provider-call accounting and deterministic reconciliation confirm model\_calls\_used = 721 (stages: shared\_initial\_generation = 720; repair = 1) and complete\_pair\_count = 720 (analysis/deterministic\_reconciliation.json, sha256 49731a2e\allowbreak{}52e049cc\allowbreak{}f5b53c35\allowbreak{}ddac2351\allowbreak{}ad75e7e9\allowbreak{}5504dc4a\allowbreak{}42bcb096\allowbreak{}d0dec3aa). Because shared\_initial\_candidate semantics were enforced, the single guarded-only detection (the unique n10 entry) must arise from downstream guarded-stage activity (sandboxed emulation) or from the single recorded repair-stage invocation.

Explicit archival pointer to the discordant episode (artifact-grounded). The unique discordant episode can be located in the archived raw-results file: execution/raw\_results.jsonl (input-results SHA-256 = df367ecc\allowbreak{}40a6f0ea\allowbreak{}f4021292\allowbreak{}60e11d84\allowbreak{}50884a4e\allowbreak{}a5fdb94a\allowbreak{}602e0d25\allowbreak{}13aef0e8). Within that file, filter for the entry where baseline\_score = 0 and guarded\_score = 1; that raw\_results.jsonl record contains the episode's pair\_id and the linked validator\_trace and any repair\_provider\_trace fields that document the repair-stage invocation. The deterministic reconciliation artifact (analysis/deterministic\_reconciliation.json, sha256 49731a2e52e0...) and the provider\_call\_audit block therein corroborate that the repair-stage call recorded in provider accounting maps to that unique discordant pair. We provide these exact artifact identifiers so readers can unambiguously locate the discordant pair and its repair/emulation traces in the archived bundle (see Reproducibility pointers below).

Discordant attribution and repair linkage. The archived raw\_results entry for the discordant pair includes validator\_trace (validation\_before/after), any emulation observations, and fields capturing whether a repair prompt was issued and its trace path. The deterministic\_reconciliation.json provider\_call\_audit documents stage\_counts and confirms exactly one repair-stage call; matching timestamps and trace-path fields in raw\_results.jsonl link that repair-stage call to the guarded-only detection. Thus the run artifacts substantiate that the single guarded-only detection corresponds to downstream guarded-stage activity (repair/emulation) rather than to generation differences, consistent with shared\_initial\_candidate semantics.

Representative per-episode diagnostics. Representative scoring traces archived under execution/scoring/task-000001-*.json ... task-000006-*.json include detailed validator\_trace objects with normalized\_configuration, validation\_before/after states, parsed\_command\_count, required\_command\_count, violation\_count, and validator\_version = 5.0. Those representative artifacts illustrate that the vast majority of shared\_initial candidates complied with validator rules (violation\_count = 0) and that sandboxed emulation generally corroborated validator findings in this curated task bank. The analysis/condition\_summary.csv artifact (sha256 394bc690671b7...) summarizes successes/failures by condition and reproduces the marginal success rates reported above.

Synthesis of numeric outcomes. In short, the guarded pipeline detected one more successful item than the baseline (720 vs. 719) across 720 paired episodes. The paired estimate (\textasciitilde{}0.00139) and its 95\% bootstrap CI ([0.0, 0.0041667]) are both orders of magnitude smaller than our preregistered minimum effect-of-interest (0.20). Provider-call accounting (model\_calls\_used = 721, repair = 1) further shows that the observed discordance is attributable to recorded downstream activity rather than to between-condition generation variance. These concrete, artifact-grounded outcomes are the basis for the interpretations and operational boundary conditions we discuss elsewhere in the manuscript.

\section{Discussion}
Why was the guarded-stage aggregate benefit negligible? The archived evidence supports three artifact-grounded factors and suggests practical boundary conditions where guarded stages are---and are not---likely to matter.

1) Validator ceiling effect and task-rule alignment. Deterministic\_netops\_validator\_v5 explicitly encodes the syntactic and semantic checks exercised by the curated tasks (examples of patterns in task payloads include shutdown\_configure\_restore, make\_before\_break\_uplink\_switchover, and paired\_link\_atomic\_mtu\_change). Representative validator\_trace fields (parsed\_command\_count, required\_command\_count, violation\_count = 0) in execution/scoring/task-000001-*.json show that most shared\_initial candidates already satisfied the validator's rule set. The baseline marginal success rate of 719/720 (0.9986111111111111) leaves essentially no headroom for the guarded emulation stage to flag additional failures in this constrained regime.

2) Curated inclusion and reduced adversarial exposure. The confirmatory sample was assembled from a deterministic generator (intent\_configuration\_repair\_v1) and required vendor-parser acceptance before inclusion. This inclusion rule intentionally restricted confirmatory items to parser-validated, policy-constrained scenarios; as a result, many runtime-only faults that emulation detects in more diverse or adversarial settings are underrepresented. The missingness summary and deterministic\_reconciliation artifacts confirm zero irrecoverable parser exclusions, indicating the set was tightly constrained to inputs the validator could reasonably analyze.

3) Shared\_initial\_candidate semantics and limited repair activity. By design we reused the exact sampled initial candidate for both baseline and guarded conditions; this isolates downstream guarded-stage signals as the only allowable source of between-condition discordance. Provider\_call\_audit and deterministic\_reconciliation report model\_calls\_used = 721 and stage\_counts: shared\_initial\_generation = 720, repair = 1. The single n10 discordant pair therefore corresponds to the single recorded repair-stage invocation or to a guarded-emulation observation not codified as a validator rule. The raw\_results and deterministic\_reconciliation artifacts together allow an auditor to locate that unique pair\_id and inspect validator\_trace and repair\_provider\_trace to see whether the repair altered the candidate or whether the emulation exposed a runtime-only violation.

Interpretation and operational implications. For environments that maintain a high-coverage static validator tailored to expected intent patterns, and where inputs are parser-validated and curated, adding an automated sandboxed-emulation stage may have negligible incremental detection benefit at scale. This does not imply sandboxing is unnecessary; rather, it implies a conditional efficiency boundary: when validator coverage closely matches expected task patterns and the input pipeline enforces strong prevalidation, guarded runtime checks will detect relatively few additional faults. Conversely, in operational regimes characterized by (a) noisy natural-language intents, (b) adversarially crafted misconfigurations, (c) broader uncurated vendor-CLI diversity, or (d) multi-step stateful interactions, the guarded stage can surface runtime semantics and side-effects that static validators cannot anticipate.

Practical recommendations grounded in the artifacts. The archived execution and analysis artifacts suggest three pragmatic paths depending on operator goals: (i) invest in richer static validators that capture more workflow-level constraints when low-latency, high-throughput intent\ensuremath{\rightarrow}config translation is the priority; (ii) deploy guarded sandboxing selectively for high-risk or high-impact changes where emulation fidelity matters; (iii) combine both approaches and use adversarial/fuzzing task banks to identify validator gaps---our evidence package includes representative traces that can seed such adversarial test-case design.

Limitations of this interpretation rooted in archived evidence. The near-zero aggregate effect we observed is specific to the preregistered task bank, the single validator implementation, the declared/requested model, and the shared\_initial\_candidate semantics. The execution manifest and deterministic\_reconciliation document the exact run accounting (model\_calls\_used = 721; repair = 1) and allow independent verification that the observed discordance is localized to a single downstream event. Because the study was powered a priori for a minimum detectable paired difference of 0.20, the empirical observation (\ensuremath{\approx}0.0014) lies far below that sensitivity and must be interpreted as a narrowly scoped near-zero finding for this experimental regime---not as broad evidence that guarded stages are always unnecessary.

Directions where guarded stages are likely to add value. The artifact-grounded diagnostics indicate three concrete extensions where guarded emulation should be retested: (1) adversarially generated tasks specifically designed to exercise validator gaps (e.g., subtle state-dependent ordering errors); (2) expanded vendor-CLI families and syntactic/semantic diversity to challenge parser and validator coverage; (3) higher-fidelity emulation that models vendor-specific control-plane subtleties and long-lived state transitions. The archived per-episode traces in execution/raw\_results.jsonl and the representative scoring artifacts provide concrete seed cases and validator-rule examples to guide these follow-up efforts.

Concluding synthesis. The run's provenance and analysis artifacts permit an unambiguous, auditable mapping from the reported aggregate counts to the single discordant episode and the recorded repair-stage call. These artifacts support a principled, evidence-led conclusion: in this preregistered, curated regime, a strong deterministic validator absorbed almost all detectable faults and the guarded sandboxed-emulation stage produced a vanishing marginal detection gain. Whether that conclusion generalizes to more adversarial or production-like regimes requires the targeted extensions described above, for which the current archived bundle provides reproducible starting points.

\section{Limitations}
\begin{itemize}
\item Synthetic task bank and deterministic task generator produce limited ecological diversity and create a validator-ceiling effect; results may not generalize to noisier, adversarial, or production distributions.
\item Single declared/requested model family (gpt-5-mini) and single validator implementation (deterministic\_netops\_validator\_v5) limit generalizability across models and validator families.
\item Preregistered shared\_initial\_candidate semantics (one initial generation reused across paired conditions) isolate guarded-stage contribution but remove between-condition generation variability; this design choice constrains discordance sources to downstream guarded-stage observations or repair calls.
\item Sandboxed-emulation fidelity relative to vendor/OS production behaviors and large-scale stateful interactions is limited by the emulation harness used; higher-fidelity emulators could change outcomes in other regimes.
\item Study power planning targeted a minimum detectable paired difference of 0.20; the observed aggregate effect (\textasciitilde{}0.0014) lies far below that design sensitivity and should be interpreted as a narrowly scoped near-zero finding for this experimental regime rather than a general negative result for all NetOps workflows.
\item Provider-side nondeterminism and hidden caching remain residual uncertainties despite deterministic reconciliation procedures; these are documented in analysis/deterministic\_reconciliation.json and provider\_call\_audit artifacts.
\end{itemize}

\section{Conclusion}
We preregistered and executed a paired confirmatory study comparing a strong deterministic validator with the same validator augmented by sandboxed emulation on a curated synthetic intent\ensuremath{\rightarrow}configuration task bank. Over 720 paired tasks the guarded pipeline produced a negligible aggregate detection gain (0.001388888888888884) with exact McNemar p = 1.0 and 95\% bootstrap CI [0.0, 0.004166666666666667]. Run accounting shows one repair-stage invocation corresponding to the single guarded-only detection; because shared\_initial\_candidate semantics were enforced, archived per-episode traces permit independent verification of that mapping via execution/raw\_results.jsonl and analysis/deterministic\_reconciliation.json. We interpret the result as arising from validator ceiling effects and curated task design; follow-up work should focus on adversarial and production-like task banks, multi-validator ensembles, and higher-fidelity emulation to identify regimes where guarded stages add substantive operational value.

Reproducibility pointers (compact). The canonical execution bundle archived by the run contains: execution/execution\_manifest.json (execution summary and preregistration SHA-256), execution/raw\_results.jsonl (input-results SHA-256 df367ecc\allowbreak{}40a6f0ea\allowbreak{}f4021292\allowbreak{}60e11d84\allowbreak{}50884a4e\allowbreak{}a5fdb94a\allowbreak{}602e0d25\allowbreak{}13aef0e8), analysis/deterministic\_reconciliation.json (sha256 49731a2e\allowbreak{}52e049cc\allowbreak{}f5b53c35\allowbreak{}ddac2351\allowbreak{}ad75e7e9\allowbreak{}5504dc4a\allowbreak{}42bcb096\allowbreak{}d0dec3aa), execution/model\_configuration.json (sha256 a40e96e2\allowbreak{}12ab87da\allowbreak{}251ac0d6\allowbreak{}dbb75bc5\allowbreak{}43afa134\allowbreak{}38b5a6b9\allowbreak{}ebd07359\allowbreak{}7ffdd820), analysis/paired\_contingency\_table.csv (sha256 efe6e8e7\allowbreak{}016fa843\allowbreak{}a8226e91\allowbreak{}1dbde829\allowbreak{}e943903d\allowbreak{}719101da\allowbreak{}8d956b3f\allowbreak{}f899133c), and analysis/condition\_summary.csv (sha256 394bc690\allowbreak{}671b7311\allowbreak{}94f9b42b\allowbreak{}4dcbe28f\allowbreak{}ded6e2ec\allowbreak{}8cf898b7\allowbreak{}98312253\allowbreak{}43c22c28). The discordant episode is the unique raw\_results.jsonl entry with baseline\_score = 0 and guarded\_score = 1; that record contains the pair\_id and validator/emulation traces that map it to the single repair-stage invocation recorded in deterministic\_reconciliation.json. These artifacts are archived as the canonical reproducibility bundle referenced by the execution manifest and the preregistration record. (See Disclosure for exact manifest references.)

\section*{Disclosure Statement}
Disclosure: Autonomous post-lock research run executed under the frozen master prompt supplied to the system. Declared/requested model: "gpt-5-mini" (provider recorded as openai\_responses). Deterministic validator: deterministic\_netops\_validator\_v5. Orchestration adapter family: hosted\_netops\_gvr\_v1. Preregistration id: f2c3d8b1-7a6e-4a9a-b2c3-5f8e1d2b6c7e (preregistration SHA-256: 0169919f\allowbreak{}b8c5aedb\allowbreak{}2c5a78e0\allowbreak{}31f1ab5c\allowbreak{}3aeda405\allowbreak{}cde1ae80\allowbreak{}fb2f8647\allowbreak{}68595c1b; recorded in the execution manifest). Initial master-prompt immutable reference: provenance/master\_prompt.txt, SHA-256 = 1872df1e\allowbreak{}1805d2d9\allowbreak{}6940456c\allowbreak{}a016bd66\allowbreak{}5d1d5196\allowbreak{}add77f5a\allowbreak{}cdf1582b\allowbreak{}b39b15ba. Execution summary (artifact-grounded): completed\_episode\_count = 1440, complete\_pair\_count = 720, model\_calls\_used = 721 (720 shared initial + 1 repair-stage call); stage\_counts and provider-call audit are archived in analysis/deterministic\_reconciliation.json (sha256 49731a2e52e0...). Human role and autonomy boundary: a human Research Director selected the broad NetOps topic family and constrained the pre-lock master prompt; after the autonomy lock the agentic pipeline autonomously performed literature selection, preregistration, code and prompt generation, experiment execution, analysis, peer review, and manuscript drafting. No manual human editing of technical manuscript content occurred after the autonomy lock. The execution manifest and archived artifacts named above serve as the canonical reproducibility bundle.

\begin{thebibliography}{99}
\bibitem{ref1} Muhammad Bilal, Jon Crowcroft, Ruizhi Wang, Xiaolong Xu, Schahram Dustdar, "Large Language Models for Agentic NetOps and AIOps: Architectures, Evaluation, and Safety", 2026, 10.48550/arxiv.2605.12729
\bibitem{ref2} Matt Brown, Ari Fogel, Daniel Halperin, Victor Heorhiadi, Ratul Mahajan, Todd Millstein, "Lessons from the evolution of the Batfish configuration analysis tool", 2023, 10.1145/3603269.3604866
\bibitem{ref3} http://arxiv.org/abs/2407.08249
\end{thebibliography}

\end{document}
